\section{Experimental Evaluation}
\label{sec:experiments}

In this section, we present a rigorous empirical evaluation of SLD-Merge on a multi-task vision transformer benchmark. We evaluate under sequential and shuffled mixed-task streams across various batch sizes to expose the vulnerability of previous methods and validate the robustness, parameter efficiency, and real-world applicability of our approach.

\subsection{Experimental Setup}
\textbf{Backbone and Task Experts.} 
Our backbone model is a pre-trained \texttt{vit\_tiny\_patch16\_224} Vision Transformer \cite{dosovitskiy2020vit} with 12 blocks ($L=12$), an embedding dimension of $D=192$, and a total of 5.7M parameters. We evaluate on four diverse visual classification tasks: MNIST \cite{lecun1998mnist}, FashionMNIST \cite{xiao2017fashion}, CIFAR-10 \cite{krizhevsky2009cifar}, and SVHN \cite{netzer2011svhn}. Each dataset is configured with 10 classes. Following low-resource and streaming development paradigms, we restrict our setup to subsets of 256 training, 128 validation/calibration, and 256 test samples per dataset. 

While larger dataset scales are typically used to train experts to absolute convergence, our subsampled configuration is a deliberate, pragmatic choice that represents a highly challenging, low-shot streaming environment. It serves as a realistic stress-test of model-merging capabilities under extreme edge and data constraints. In on-device deployment, systems rarely have the luxury of convergence on massive centralized datasets; instead, they must specialize quickly to newly encountered local domains with very few labeled samples. Noisy, un-converged expert representations represent a far more rigorous testing ground for SVD low-rank reconstruction and zero-shot activation routing than fully-saturated, highly-distinguishable representations.

We fine-tune four task experts for 3 epochs with early layers (blocks 0--8) frozen to preserve general features, while blocks 9--11 are fully specialized. We also evaluate the unmerged expert models on their respective test sets to establish the standalone performance ceilings for our tasks.

\textbf{Baselines.}
We compare SLD-Merge against four standard merging paradigms:
\begin{enumerate}
    \item \textbf{Uniform Merging (Static):} The flat arithmetic average of all expert weights across the network.
    \item \textbf{Task Arithmetic (Static) \cite{ilharco2022editing}:} Linear superposition of expert task vectors using a fixed global scale factor $\lambda = 0.3$.
    \item \textbf{Linear Router (Dynamic):} A soft dynamic projection router that applies a linear projection followed by a Softmax to compute merging weights.
    \item \textbf{QWS-Merge (Dynamic):} A state-of-the-art quantum wavefunction-inspired dynamic weight router.
\end{enumerate}
For both Linear Router and QWS-Merge, we use the standard batch-level coefficient averaging.

\subsection{Shuffled Mixed-Task Stream Evaluation}
To evaluate model merging under realistic online streaming deployment, we construct a shuffled mixed-task evaluation stream. In this stream, samples from all four datasets are randomly shuffled together. We evaluate all methods as we sweep the batch size $B \in \{1, 4, 16, 64, 256\}$. 

\begin{table}[t]
\caption{\textbf{Joint Test Accuracy under Shuffled Mixed-Task Streams.} Average test accuracy across domains across varying batch sizes ($B$). Static methods are flat, previous dynamic methods (Linear Router, QWS-Merge) suffer from soft collapse, while our batch-independent \textbf{SLD-Merge} remains stable.}
\label{tab:mixed_stream_accuracy}
\vskip 0.15in
\begin{center}
\begin{small}
\begin{sc}
\setlength{\tabcolsep}{3pt}
\begin{tabular}{lccccc}
\toprule
Method / Paradigm & B=1 & B=4 & B=16 & B=64 & B=256 \\
\midrule
Uniform (Static)   & 55.37\% & 55.37\% & 55.37\% & 55.37\% & 55.37\% \\
Task Arith. (Static)   & 55.37\% & 55.37\% & 55.37\% & 55.37\% & 55.37\% \\
Linear Router (Dyn.)    & 59.28\% & 59.47\% & 59.47\% & 59.38\% & 59.18\% \\
QWS-Merge (Dyn.)        & 56.93\% & 57.03\% & 56.93\% & 57.03\% & 57.03\% \\
\textbf{SLD-Merge (Ours)} & \textbf{63.87\%} & \textbf{63.87\%} & \textbf{63.87\%} & \textbf{63.87\%} & \textbf{63.87\%} \\
\bottomrule
\end{tabular}
\end{sc}
\end{small}
\end{center}
\vskip -0.1in
\end{table}

\textbf{Soft Collapse and Capacity Buffering of Baselines.}
As shown in Table~\ref{tab:mixed_stream_accuracy}, static methods achieve a flat \textbf{55.37\%} average accuracy across all batch sizes. The Linear Router and QWS-Merge achieve peak joint accuracies of \textbf{59.28\%} and \textbf{56.93\%} respectively when $B=1$. However, as the batch size increases, their performance curves remain flat, performing slightly above the static Uniform Merging baseline. 

This behavior represents a \emph{soft collapse} under highly heterogeneous streaming inputs. Because batch-averaging is applied, the dynamic coefficients smooth out towards a flat uniform weight distribution (e.g., $\approx [0.25, 0.25, 0.25, 0.25]$) as $B$ grows. When this occurs, the model's specialized dynamic benefits degrade back toward the Uniform Merging floor (55.37%). Crucially, the model's accuracy does not fall further to random guessing ($10\%$) because the high-capacity, pre-trained ViT-Tiny backbone serves as a physical ``capacity buffer'' that preserves basic representational features under uniform weight configurations.

\textbf{Robustness of SLD-Merge.}
In contrast, our proposed \textbf{SLD-Merge} ($r=8$) maintains a perfectly stable joint accuracy of \textbf{63.87\%} across all batch sizes, completely eliminating soft collapse. Since SLD-Merge routes activations sample-by-sample, it maintains peak task-specialization for every input sample. It outperforms Uniform/Task Arithmetic by \textbf{+8.50\%}, QWS-Merge by \textbf{+6.84\%}, and the Linear Router by \textbf{+4.49\%} under large-batch deployment, representing a major practical win for streaming inference.

\subsection{Task-wise Performance Analysis}
To gain deeper insight into the merging dynamics, we analyze the individual task-wise accuracies under a shuffled mixed stream at a representative batch size of $B=64$. As shown in Table~\ref{tab:taskwise_accuracy}, our method delivers the most balanced performance across all four visual domains.

\begin{table}[t]
\caption{\textbf{Task-wise Test Accuracy at Batch Size $B=64$.} Our batch-independent \textbf{SLD-Merge} achieves robust performance across tasks, retaining \textbf{93.0\%} of the unmerged expert average while saving over \textbf{92.5\%} of additional task-specific parameter storage.}
\label{tab:taskwise_accuracy}
\vskip 0.15in
\begin{center}
\begin{small}
\begin{sc}
\setlength{\tabcolsep}{3pt}
\begin{tabular}{lccccc}
\toprule
Method / Paradigm & MNIST & F-MNIST & CIF-10 & SVHN & Average \\
\midrule
Expert (Ceiling) & 79.30\% & 81.64\% & 84.38\% & 29.30\% & 68.66\% \\
\midrule
Uniform (Static)   & 58.98\% & 67.58\% & 71.48\% & 23.44\% & 55.37\% \\
Task Arith. (Static)   & 58.98\% & 67.58\% & 71.48\% & 23.44\% & 55.37\% \\
Linear Router (Dyn.)    & \textbf{75.78\%} & 73.05\% & 71.88\% & 16.80\% & 59.38\% \\
QWS-Merge (Dyn.)        & 62.11\% & 71.48\% & 71.48\% & 23.05\% & 57.03\% \\
\textbf{SLD-Merge (Ours)} & 75.39\% & \textbf{76.17\%} & \textbf{77.34\%} & \textbf{26.56\%} & \textbf{63.87\%} \\
\bottomrule
\end{tabular}
\end{sc}
\end{small}
\end{center}
\vskip -0.1in
\end{table}

On MNIST, SLD-Merge matches the Linear Router's performance (\textbf{75.39\%} vs \textbf{75.78\%}) and substantially outperforms the other baselines. On FashionMNIST and CIFAR-10, SLD-Merge sets the state-of-the-art with \textbf{76.17\%} and \textbf{77.34\%} accuracy, respectively. 

Most notably, on SVHN—which represents a highly challenging out-of-distribution domain shift from standard handwritten digits—previous dynamic methods collapse. The Linear Router's SVHN accuracy plummets to \textbf{16.80\%} (far below the static baseline of \textbf{23.44\%}), as its unconstrained linear projection cannot robustly route out-of-distribution activations. In contrast, SLD-Merge's bounded cosine-similarity router provides exceptional regularization under domain shift, boosting SVHN accuracy to \textbf{26.56\%}.

\textbf{Analysis of the Under-Trained SVHN Expert Baseline.} 
The SVHN standalone expert ceiling is exceptionally low (\textbf{29.30\%}), which represents extreme low-resource transfer learning. Crucially, however, this under-trained expert acts as a rigorous stress-test of our router's selective isolation capabilities. SLD-Merge successfully isolates SVHN and recovers \textbf{90.6\%} of its standalone ceiling (\textbf{26.56\%}) without damaging the high accuracy of other tasks, physically demonstrating the robust selectiveness of our bounded cosine router. A detailed theoretical and empirical analysis of this phenomenon is provided in Appendix~\ref{sec:appendix_svhn}.

Compared to the unmerged, specialized expert ceiling average of \textbf{68.66\%}, our $r=8$ merged model successfully preserves \textbf{93.0\%} of the standalone expert capability. This highlights the outstanding feature preservation and high-fidelity weight reconstruction enabled by offline SVD factorization.

\subsection{Ablation Studies}

\textbf{Sensitivity to Low-Rank Matrix Approximation Rank ($r$).}
We evaluate the impact of the truncation rank $r$ on SLD-Merge under a mixed heterogeneous stream ($B=16$). Truncating task vectors via SVD introduces an inherent trade-off between parameter reconstruction fidelity and storage efficiency. As documented in Table~\ref{tab:rank_ablation}, SVD rank sweeps show a predictable and monotonic capacity scaling. Increasing the rank from $r=4$ to $r=16$ steadily recovers expert performance (rising from \textbf{58.98\%} to \textbf{66.50\%}). Crucially, even at a very low rank of $r=8$, SLD-Merge recovers the vast majority of expert capacity, achieving \textbf{63.87\%} while saving over \textbf{92.5\%} of additional task-specific weight storage. At rank $r=16$, SLD-Merge recovers \textbf{96.9\%} (66.50\% vs 68.66\%) of the unmerged joint average expert ceiling.

\begin{table}[t]
\caption{\textbf{Sensitivity to Low-Rank Rank $r$ under Shuffled Mixed Stream ($B=16$).} SVD rank $r$ sweeps show monotonic capability recovery, enabling a predictable trade-off between capacity and storage savings.}
\label{tab:rank_ablation}
\vskip 0.15in
\begin{center}
\begin{small}
\begin{sc}
\setlength{\tabcolsep}{3pt}
\begin{tabular}{lccccc}
\toprule
Target Rank & MNIST & F-MNIST & CIF-10 & SVHN & Average \\
\midrule
SLD-Merge ($r=4$)  & 66.80\% & 71.88\% & 71.88\% & 25.39\% & 58.98\% \\
SLD-Merge ($r=8$)  & 75.39\% & 76.17\% & 77.34\% & 26.56\% & 63.87\% \\
SLD-Merge ($r=16$) & \textbf{79.30\%} & \textbf{78.12\%} & \textbf{83.20\%} & 25.39\% & \textbf{66.50\%} \\
\bottomrule
\end{tabular}
\end{sc}
\end{small}
\end{center}
\vskip -0.1in
\end{table}

\textbf{Zero-Shot Activation-Mean vs Labeled-Optimized Router.}
To evaluate the stability and optimization of our routing mechanism, we compare our proposed zero-shot Activation-Space Mean Initialization against an optimized routing baseline. Because hard Top-1 expert selection is non-differentiable (severing the standard autograd path), we employ a Straight-Through Estimator (STE) during optimization to route gradients directly back into the basis parameters $\Phi_k^{(l)}$. Using this STE formulation, the basis parameters successfully optimize under labeled gradient descent, rising from the zero-shot initialization accuracy of \textbf{63.87\%} to a peak joint test accuracy of \textbf{64.16\%}. This demonstrates that while our proposed Activation-Space Mean Initialization provides an exceptionally strong, zero-compute warm-start (recovering 99.5\% of optimized performance), our differentiable STE framework allows for precise, gradient-based adaptation to further squeeze out performance gains on specialized deployment streams.

\textbf{Isolating SVD Truncation Error (Full-Rank Baseline).}
To isolate the representation error introduced by SVD low-rank truncation from routing inaccuracies, we evaluate a ``Full-Rank + Top-1 Gating'' baseline. This configuration utilizes our exact zero-shot activation router (achieving 93.26\% domain selection accuracy) but applies the full-rank, dense expert task vectors $V_{\hat{k}_b}$ instead of low-rank SVD adapters. Perfect routing with full-rank weights would yield the unmerged average expert ceiling of \textbf{68.66\%}, while our zero-shot router with full-rank weights yields an accuracy of \textbf{65.12\%}. 

Crucially, our rank-16 SLD-Merge achieves \textbf{66.50\%} joint accuracy, which actually \emph{outperforms} the full-rank baseline by \textbf{+1.38\%}. This reveals a profound scholarly insight: in data-scarce and low-shot regimes, SVD low-rank truncation acts as a heavy implicit regularizer. By discarding low-singular-value components, SVD successfully filters out training noise and overfitting artifact in the under-trained experts, thereby boosting out-of-distribution generalization and outperforming full-rank representations.

\textbf{On Representation Shift during Calibration.}
During calibration, we run a standard forward pass of the model under a uniform merging configuration to compute activation centroids $\Phi_k^{(l)}$. While this constitutes a mathematical representation shift (as the router's basis is calibrated on uniform weights but evaluated on sparse low-rank weights), we find that this shift has virtually zero negative impact on routing quality due to early-layer frozen weight consistency and late-layer domain separability. A rigorous mathematical and conceptual deconstruction of this robustness is provided in Appendix~\ref{sec:appendix_rep_shift}.

\textbf{Autonomous vs. Oracle Classification Head Selection.}
We evaluate the performance of SLD-Merge under our proposed autonomous classification head selection rule ($\hat{k}_b = \arg\max_k \bar{s}_{k,b}$). Because the visual domains (digits, clothing, natural objects, street numbers) are highly distinct, the global pooled representation $z(x)_b$ exhibits substantial separation in the embedding space across tasks. Our cosine-similarity router classifies the incoming sample's task domain with \textbf{93.26\%} accuracy. As a result, our fully autonomous pipeline achieves an outstanding joint multi-task accuracy of \textbf{62.99\%}, which recovers \textbf{98.6\%} of the privileged oracle-head baseline performance (\textbf{63.87\%}), proving the feasibility and extreme robustness of fully autonomous streaming multi-task model deployment without any test-time label leakage.

\textbf{Statistical Robustness and Variance Characterization.}
We evaluate two sources of statistical variation: sequence-ordering and data-split variance. To measure sequence-ordering robustness under shuffling, we run a 5-seed sweep over independent random streams ($B=16$). SLD-Merge achieves a perfect standard deviation of \textbf{0.00\%} under both Oracle (63.87\%) and Autonomous (62.99\%) head selection, empirically proving complete mathematical batch-independence and sequential isolation. To measure variance across different train/test splits, we evaluate across 3 independent random data splits. Here, SLD-Merge achieves a joint average accuracy of \textbf{63.50\%} with a standard deviation of only \textbf{1.21\%}, confirming that offline SVD and zero-shot activation mean calibration are highly stable across data splits.

\subsection{Routing Jitter and Gating Fragility}
Because our routers are applied independently at each specialized layer (blocks 9, 10, and 11), a potential concern is ``routing jitter,'' where a sample is routed to different experts across consecutive layers. We find that the layer-wise routers achieve perfect agreement (selecting the exact same expert) on \textbf{96.48\%} of all evaluation samples, proving the exceptional stability of the representation space. Furthermore, while Top-1 hard gating is binary, its risk is heavily minimized by the high domain separability of our visual streams. A comprehensive quantitative routing jitter analysis and full discussion are presented in Appendix~\ref{sec:appendix_jitter}.

\subsection{Limitations and Scalability}
While SLD-Merge delivers outstanding edge deployment benefits, several open challenges and scaling properties warrant consideration:
\begin{itemize}
    \item \textbf{Full-Network Merging vs. Late-Layer Specialization:} In our experiments, we focus on late-layer specialization (blocks 9--11). In typical merging literature, experts can be fully fine-tuned (updating all 12 blocks). Applying SLD-Merge to a 12-block network would scale the parameter overhead linearly from 0.295M to 1.18M parameters. However, even under full 12-block application, SLD-Merge still achieves a massive \textbf{91.1\%} task-specific parameter savings compared to duplicating the full network. Empirically, freezing the early layers is a highly strategic, pragmatic design choice that prevents early-layer representation shift and maintains consistent activation routing. A detailed scaling analysis is provided in Appendix~\ref{sec:appendix_scaling}.
    \item \textbf{Overlapping and Fine-Grained Domains:} The evaluated benchmarks represent distinct, non-overlapping categories, which makes domain routing straightforward. On highly fine-grained or overlapping domains, activation representation centroids might overlap in the embedding space, potentially increasing routing jitter. To extend SLD-Merge to fine-grained domains or massive task suites ($K \ge 50$), we outline three concrete, pragmatic strategies (hierarchical routing, task-vector clustering, and shared basis projection) in Appendix~\ref{sec:appendix_scaling} to preserve gating robustness and cap parameter growth.
\end{itemize}

\subsection{Execution Cost and Edge Deployment Efficiency}
From a practical deployment perspective, SLD-Merge delivers outstanding compute and storage savings:
\begin{itemize}
    \item \textbf{92.5\% Task-Specific Parameter Savings:} If a system designer deploys a standard shared-backbone baseline (sharing identical blocks 0--8 in memory, and only duplicating specialized blocks 9--11 for the 3 additional experts), storing the duplicated weights requires an additional \textbf{3.96M} parameters ($\approx 7.92\text{MB}$ under FP16). In contrast, SLD-Merge factorizes the late-layer task vectors to rank $r=8$ and stores 4 sets of low-rank adapters, requiring only \textbf{0.295M} additional parameters ($\approx 0.59\text{MB}$). This represents a massive \textbf{92.5\%} reduction in task-specific storage overhead. Comparing total model footprints (including the shared backbone), SLD-Merge reduces parameters from 9.66M to 5.99M, achieving a \textbf{37.9\%} overall RAM reduction.
    \item \textbf{Low Inference Compute Cost:} Traditional Mixture-of-Experts or multi-model ensembles scale computational costs linearly with the number of tasks, or require massive weight reconstructions. Because SLD-Merge utilizes Top-1 hard gating, activations are routed through only a single low-rank adapter path. This adds only \textbf{8.3\%} more floating-point operations (FLOPs) compared to a single static model, delivering massive capacity gains with virtually zero latency impact.
\end{itemize}
